\section{\label{sec:wc}Wilson coefficients for the sOPE}

The procedure for calculating smeared Wilson coefficients parallels that 
for the OPE, discussed in, for example, Ref.~\cite{Collins:1984rdg}. With the 
Feynman rules of 
Sec.~\ref{sec:phi4} in hand, the calculation of smeared Wilson 
coefficients up to next-to-leading-order in perturbation theory is 
straightforward. We 
determine the smeared Wilson coefficients for the leading connected and 
disconnected operators, starting with the disconnected contribution.

\subsection{Disconnected contributions}

\subsubsection{Leading order}

We illustrate the leading-order (tree-level) and next-to-leading-order 
(one-loop) contributions to the smeared Wilson 
coefficient for the disconnected operator, $\mathbb{I}$, in Fig.~
\ref{fig:sope2pt}.
\begin{figure}[htbp]\begin{minipage}{0.49\textwidth}
\includegraphics[width=0.95\textwidth,keepaspectratio=true]{images/d1tree}
\end{minipage}%
\begin{minipage}{0.49\textwidth}
\includegraphics[width=0.95\textwidth,keepaspectratio=true]{images/d1loop}
\end{minipage}
\caption{\label{fig:sope2pt}Diagrams representing the contributions to the 
Wilson coefficient $d_{\mathbb{I}}$ at leading-order (left two 
diagrams) and next-to-leading-order (right two diagrams). Solid black and 
dashed blue 
lines are propagators at vanishing and nonvanishing flow times, respectively. 
The black squares are unsmeared fields $\phi(0)$, black dots are 
interaction vertices at vanishing flow time and the blue blob represents 
the smeared operator $\rho^2(\tau,0)$.}
\end{figure}
As we will see, the leading-order 
contribution is independent of the renormalization scale, $\mu$, because at 
this order the flow time serves as the ultraviolet regulator. Beyond leading 
order, however, the smeared Wilson coefficient will have a renormalization 
scale dependence.

We extract the disconnected smeared Wilson coefficient by considering matrix 
elements of each of the operators in the sOPE in the vacuum, which removes any 
connected contributions.  To ${\cal O}(x)$ we have
\begin{align}
\big\langle \Omega  |  {} &{\cal T}\{ 
\phi(x) \phi(0)\}  | \Omega \big\rangle =  
\frac{d_{\mathbb{I}}(\mu x,\mu^2\tau,mx)}{4\pi^2x^2}\big\langle \Omega | \,
\mathbb{I} \,| \Omega \big\rangle
\nonumber\\
{} & + 
d_{\rho^2}(\mu x,\mu^2\tau,mx)\big\langle \Omega | [\rho^2(\tau,0) 
]_{\mathrm{R}}| \Omega 
\big\rangle  + 
{\cal O}(x),
\end{align}
Here we have chosen the normalization of $d_{\mathbb{I}}$ so that at leading 
order in the free theory $d_{\mathbb{I}}$ is unity. We expand each quantity in 
this expression to one loop according to
\begin{equation}\label{eq:drhoexp}
f = f^{(0)} - \lambda f^{(1)} + {\cal O}(\lambda^2),
\end{equation}
where $f$ stands for either a matrix element or Wilson coefficient. We have 
chosen the sign of the one-loop contribution to factor out the sign of the 
coupling constant arising from the Feynman rule for the four-point vertex 
$V^{(4)} = -\lambda/24$. 

The tree-level Wilson coefficient is then given by the small spacetime 
behavior of
\begin{align}
\frac{d_{\mathbb{I}}^{(0)}}{4\pi^2x^2}  \stackrel{x\sim 0}{=}{} & 
 \Big\{\big\langle 
\Omega | 
{\cal T}\{ 
\phi(x){} \phi(0)\} | \Omega \big\rangle \nonumber\\
{} & \qquad - \big\langle \Omega | 
\rho^2(\tau,0)  | \Omega 
\big\rangle\Big\}_{{\cal O}(m^2)}^{(0)} ,
\end{align}
where we have neglected the arguments of the coefficient for clarity and used 
the fact that the tree-level connected
coefficient is $d_{\rho^2}^{(0)}=1$. The subscript indicates 
that we must expand the result to ${\cal O}(m^2)$. For more details about the 
calculation of Wilson coefficients see, for example, 
Ref.~\cite{Collins:1984rdg}. 
The corresponding Feynman integral representation is
\begin{align}
d_{\mathbb{I}}^{(0)} \stackrel{x\sim 0}{=} {} & 4\pi^2 x^2\left\{\int_k 
\frac{e^{ik 
\cdot 
x}-e^{-2k^2 \tau 
}}{k^2+m^2}\right\}_{{\cal O}(m^2)}
\nonumber \\
= {} & 4\pi^2x^2\int_k \left(e^{ik \cdot x} - e^{-2k^2 
\tau}\right)\left(\frac{1}{k^2} - \frac{m^2}{(k^2)^2}\right),
\end{align}
where
\begin{equation}
\int_k \equiv \int\frac{\mathrm{d}^4k}{(2\pi)^4}.
\end{equation}

The smeared Wilson 
coefficient is
\begin{equation}\label{eq:d1lead}
d_{\mathbb{I}}^{(0)} = 1 -
\frac{x^2}{8\tau}  +\frac{m^2x^2}{4}\left[\gamma_{\;\mathrm{E}} -1 + 
\log\left(\frac{x^2}{8\tau}\right)\right],
\end{equation}
with $\gamma_{\;\mathrm{E}}\simeq 0.577216$ the Euler--Mascheroni constant.

We can compare this result with the Wilson coefficient 
for the OPE for $d = 4-2\epsilon$ 
dimensions \cite{Collins:1984rdg}:
\begin{equation}
\overline{c}_{\mathbb{I}}^{(0)} =  
1 + \frac{  
m^2x^2}{4} \left[\frac{1}{\epsilon} + 1 +\gamma_{\;\mathrm{E}} + 
\log\left(\frac{\pi\mu^2x^2}{4}\right)\right].
\end{equation}
Then in the $\ms$ scheme, this becomes
\begin{equation}
\overline{c}_{\mathbb{I}}^{(0)} =  
1 +\frac{  
m^2x^2}{4} \left[1 +  2\gamma_{\;\mathrm{E}} + 
\log\left(\frac{\mu^2x^2}{16}\right)\right].
\end{equation}

Although the finite contribution to these expressions cannot be directly 
compared, 
because we have expressed the Wilson coefficients in two different 
renormalization schemes, we note three important features. First, the 
logarithmic dependence on the spacetime separation is identical. Second, 
the flow time $\tau$ plays the role of the renormalization scale at leading 
order. Third, we see that for small spacetime 
separations, we require a small flow-time parameter. If we do 
not choose the flow-time parameter appropriately, we generate large logarithmic 
contributions to the smeared Wilson coefficients and the sOPE exhibits poor 
convergence properties, even at small spacetime separations.

\subsubsection{Next-to-leading-order}

At one loop, the smeared Wilson coefficient is given by
\begin{align}
-\lambda\frac{d_{\mathbb{I}}^{(1)}}{4\pi^2x^2} = {} & \Big\{\big\langle 
\Omega | 
{\cal T}\{ 
\phi(x){} \phi(0)\} | \Omega \big\rangle \nonumber\\
{} & \qquad \qquad  - \big\langle \Omega | 
\rho^2(\tau,0)  | \Omega 
\big\rangle\Big\}_{{\cal O}(m^2)}^{(1)}.
\end{align}

The four-point interaction in this diagram, which we show in Fig.~
\ref{fig:sope2pt}, appears at zero flow time. Therefore the flow 
time cannot act as a regulator for the momentum integral over $k_2$ and we must 
introduce a renormalization procedure. We use dimensional regularization and 
the $\ms$ scheme. The double integral is straightforward, however, because the 
two integrals can be carried out separately:
\begin{equation}\label{eq:d11loop}
\frac{d_{\mathbb{I}}^{(1)}}{4\pi^2x^2}=  \bigg\{\int_{k_1}
\frac{e^{ik_1 \cdot 
x} - e^{- 2k_1^2\tau}}{(k_1^2+m^2)^2}  
\frac{1}{2}\int_{k_2}\frac{1}{k_2^2+m^2}\bigg\}_{{
 \cal O}(m^2)}.
\end{equation}
We find, for $m>0$,
\begin{align}
d_{\mathbb{I}}^{(1)} = {} & -
\frac{m^2x^2}{128\pi^2}\left[\gamma_{\;\mathrm{E}}-1+\log\left(\frac{x^2}{
8\tau } \right)\right]\left[
1+ \log\left(\frac{\mu^2}{m^2}\right)\right].
\end{align}

Here we see that the second term, which is a function of the new 
renormalization scale, $\mu$, and the bare mass, $m$, is nothing other than the 
one-loop contribution to the mass renormalization of the original theory, in 
the $\overline{MS}$ scheme \cite{Kleinert:2001ax}. Moreover, the factor 
containing the flow time 
is identical to the ${\cal O}(m^2)$ term from the leading-order contribution, 
$d_{\mathbb{I}}^{(0)}$, in Eq.~\eqref{eq:d1lead}.

Therefore, we can simply combine the leading-order and next-to-leading-order 
terms to give
\begin{align}
d_{\mathbb{I}} = {} & 1 -
\frac{x^2}{8\tau}  +\frac{m_{\mathrm{R}}^2x^2}{4}\left[\gamma_{\;\mathrm{E}} -1 
+ 
\log\left(\frac{x^2}{8\tau}\right)\right]  +{\cal O}(\lambda^2),
\end{align}
where $m_{\mathrm{R}}$ is the renormalized mass in the 
$\overline{MS}$ scheme, given by $m^2_{\mathrm{R}} = Z_m^{-1}m^2$ and $Z_m$ is 
the mass renormalization \cite{Kleinert:2001ax}:
\begin{equation}
Z_m = 1+\frac{\lambda}{16\pi^2}\left[
1+ \log\left(\frac{\mu^2}{m^2}\right)\right]+{\cal O}(\lambda^2).
\end{equation}

This is a clear, next-to-leading-order example of how the 
divergences of the theory at nonzero flow time are absorbed by 
the renormalization parameters of the original theory at zero flow time. In 
other words, the renormalized theory at zero flow time remains ultraviolet 
finite at nonzero flow time.

\subsection{Connected contribution}

We illustrate the leading- and next-to-leading-order contributions
to the connected Wilson coefficient $d_{\rho^2}(\mu x,\mu^2\tau,mx)$ of 
Eq.~\eqref{eq:sope2pt} in Figs.~\ref{fig:sopeconn} and 
\ref{fig:sopehigher}, respectively. 
\begin{figure}[htbp]\includegraphics[width=0.42\textwidth,keepaspectratio=true]{images/dphi1loop}
\caption{\label{fig:sopeconn} 
Diagrams representing the leading-order (one-loop) contributions to 
the Wilson coefficient $d_{\rho^2}$. Details of the 
Feynman diagrams provided in the caption of Fig.~\ref{fig:sope2pt}.}
\end{figure}
\begin{figure}[htbp]\begin{minipage}{0.6\textwidth}
\includegraphics[width=0.95\textwidth,keepaspectratio=true]{images/dphi2loopa}
\end{minipage}%
\caption{\label{fig:sopehigher} 
Diagrams representing the contributions to the next-to-leading-order 
(two-loop) contributions to the 
Wilson coefficients $d_{\rho^2}$. For details of the Feynman 
diagrams, see the caption of Fig.~\ref{fig:sope2pt}.}
\end{figure}
Throughout this section, we neglect
diagrams that are trivially incorporated as part of the wave function 
renormalization of the external fields, such as the one-loop examples 
illustrated in Fig.~\ref{fig:zphi}. Provided the original 
theory at zero flow time is renormalized, counterterms cancel these 
contributions completely.
\begin{figure}[htbp]\includegraphics[width=0.62\textwidth,keepaspectratio=true]{images/dphizphi}
\caption{\label{fig:zphi}
Example diagrams that are naturally 
incorporated in the renormalized external propagators. We do not include these 
contributions in our explicit calculations of the Wilson coefficients. For 
details of the Feynman 
diagrams, see the caption of Fig.~\ref{fig:sope2pt}.}
\end{figure}

\subsubsection{Leading order}

We extract the leading-order connected Wilson coefficient by considering matrix 
elements of each of 
the operators in the sOPE coupled to two external fields, which removes any 
disconnected contributions.
We can then read off the one-loop contribution to the Wilson coefficient, 
$d_{\rho^2}^{(1)}$, by matching terms at ${\cal O}(\lambda)$. This 
contribution is given by the small spacetime behavior of
\begin{align}
-\lambda d_{\rho^2}^{(1)} {} & \stackrel{x\sim 0}{=}  
\Big\{ \big\langle \Omega | {\cal T}\{ 
\phi(x){} \phi(0)\}  
\widetilde{\phi}(p_1) \widetilde{\phi}(p_2) | \Omega \big\rangle^{(1)} 
\nonumber\\
{} & - \big\langle \Omega | \rho^2(\tau,0) 
\widetilde{\phi}(p_1) \widetilde{\phi}(p_2) | \Omega 
\big\rangle^{(1)}\Big\}_{{\cal O}(m^0)}.
\end{align}  

The corresponding Feynman integral 
is
\begin{align}\label{eq:curlyb}
\frac{1}{(p^2_1+m^2)(p_2^2+m^2)}\left\{\frac{1}{2}\int_k 
\frac{e^{ik\cdot x} - e^{-(k^2+q^2) \tau }}{(k^2+m^2)(q^2+m^2)}\right\},
\end{align}
where $q = k-p_1-p_2$. We extract the smeared Wilson 
coefficient by examining the small spacetime behavior of the 
integral in curly braces, expanded to ${\cal O}(m^0)$:
\begin{align}\label{eq:drho1}
d_{\rho^2}^{(1)} \stackrel{x\sim 0}{=} \left\{
\frac{1}{2}\int_k 
\frac{e^{ik\cdot x} - e^{-(k^2+q^2)\tau }}{(k^2+m^2)(q^2+m^2)}\right\}_{{\cal 
O}(m^0)}.
\end{align}

The 
smeared Wilson coefficients must be 
independent of the external states to ensure that the sOPE is truly an operator 
expansion. In this particular case, 
we require that 
$d_{\rho^2}^{(1)}$ is independent of the external momenta 
$p_1$ and $p_2$ and the mass. By taking 
a derivative with respect to one of the external momenta,
\begin{equation}\label{eq:drhoderiv}
\frac{\mathrm{d}}{\mathrm{d}p_i}d_{\rho^2}^{(1)} =
 \int_k \frac{q_i\big[e^{ik\cdot x} - 
e^{-(k^2+q^2)\tau}(1+(q^2+m^2)\tau)\big]}{(k^2+m^2)(q^2+m^2)^2},
\end{equation}
we obtain a convergent integral. The $x\rightarrow 0$ limit of this integral is 
now well defined and only vanishes if the flow time is 
related to the spacetime separation. An analogous result holds if we take a 
derivative with respect to the external mass, $m$. On dimensional grounds, 
then, we 
guarantee that the smeared Wilson coefficient is independent of the external 
states by choosing $\tau \propto x^2$.

On physical 
grounds, however, we require that the smearing radius is 
smaller than the spacetime extent of the nonlocal operator: 
$s_{\mathrm{rms}}< x$. This choice ensures that the sOPE remains an expansion 
in local operators. Physically speaking, if the 
gradient flow probes length scales on the order of 
the nonlocal operator, which we represent in Fig.~\ref{fig:sopeexpansion}, 
then smeared operators would cease to be (approximately) local. In other words, 
the 
sOPE would be a poor expansion for the original operator. This is the 
physical origin of the third feature that we saw in the leading-order 
disconnected contribution, Eq.~\eqref{eq:d1lead}: small spacetime 
separations require small flow times to ensure convergence. The OPE requires 
small spacetime 
separations for good convergence, 
so it follows that the sOPE requires small flow times as well. 
\begin{figure}[htbp]\includegraphics[width=0.62\textwidth,keepaspectratio=true]{images/dphiexpansion}
\caption{\label{fig:sopeexpansion} At small flow times (left-hand diagram), 
the smeared operators are localized relative to the spacetime separation of 
the nonlocal operator. At large flow-time values (right-hand diagram), 
however, 
the smearing radius probes the scale of the nonlocal and the sOPE is a poor 
approximation to the original nonlocal operator. For details of the Feynman 
diagrams, see the caption of Fig.~\ref{fig:sope2pt}.
}
\end{figure}

There is a complementary viewpoint that elucidates the small flow-time 
requirement more quantitatively: the small flow-time expansion 
\cite{Makino:2014taa,Suzuki:2013gza,Luscher:2011bx}. In this case we see that 
the 
derivative of Eq.~\eqref{eq:drhoderiv} is independent of the external 
momenta in the $x\rightarrow 0$ limit, up to terms linear in the flow time:
\begin{equation}
\frac{\mathrm{d}}{\mathrm{d}p_i}d_{\rho^2}^{(1)} =
\int_k\frac{q_i\big(e^{ik\cdot x} - 
1\big)}{(k^2+m^2)(q^2+m^2)^2}+{\cal O}(\tau).
\end{equation}
From this, it is clear that the Wilson coefficients will be independent of the 
external states, up to terms linear in the flow time. Therefore, provided the 
flow time is small relative to the spacetime separation, the sOPE is an 
expansion in approximately local operators, in a quantifiable sense. Moreover, 
the 
flow-time dependence will cancel in the product of the Wilson coefficient and 
its associated matrix element, to the desired order in the flow time.

We are free to set $p_1 = p_2 = 0$ in Eq.~
\eqref{eq:drho1}, because the smeared Wilson coefficient 
is independent of the 
external momenta and the mass to the order at which we are 
working \cite{Collins:1984rdg}. Expanding in the mass, we obtain
\begin{align}
d_{\rho^2}^{(1)} = {} &  
\frac{1}{2}\int_k 
\frac{e^{ik\cdot x} - e^{-2k^2 \tau }}{(k^2)^2} \nonumber \\
 = {} & -\frac{1}{32\pi^2}\left[\gamma_{\;\mathrm{E}} -1 + 
\log\left(\frac{x^2}{8\tau}\right)\right].
\end{align}
Combining this with the leading-order contribution, which is just unity, we have
\begin{equation}\label{eq:drhores}
d_{\rho^2}= 1 +\frac{\lambda}{32\pi^2}\left[\gamma_{\;\mathrm{E}} -1 + 
\log\left(\frac{x^2}{8\tau}\right)\right]+{\cal O}(\lambda^2).
\end{equation}

In contrast, the Wilson coefficient for the OPE in the $\ms$ scheme, denoted by 
$\overline{c}$, is
\begin{equation}\label{eq:cphires}
\overline{c}_{\phi^2} =  1+\frac{\lambda}{32\pi^2} \left[1 
+2\gamma_{\;\mathrm{E}}+ 
\log\left(\frac{\mu^2x^2}{16}\right)\right]+{\cal O}(\lambda^2).
\end{equation}
We note the occurrence of the three features we observed for the leading-order 
disconnected contribution: the same spacetime dependence for both smeared and 
unsmeared coefficients; the appearance of the flow time as a leading-order 
regulator; and the need to choose a small flow time for small spacetime 
separations to avoid large logarithmic contributions. From the small flow-time 
expansion viewpoint, we can confirm that the flow-time 
dependence ultimately cancels to the desired order. For example, the matrix 
element of $\rho^2(\tau,0)$ coupled to two external fields is
\begin{align}
\!\!\!\!\big\langle \Omega 
|{} & \rho^2(\tau,0)\widetilde{\phi}(p_1) \widetilde{\phi}(p_2)
|\Omega \big\rangle = 1  \nonumber\\
{} &\quad +\frac{\lambda}{32\pi^2}\left[1+\gamma_{\;\mathrm{E}} + 
\log\left(2m^2\tau\right)\right]
+{\cal O}(\tau,\lambda^2)
\end{align}
for sufficiently small flow times. Here we have dropped the external fields, 
which we take to be on shell, for simplicity. The product of this 
matrix element with the Wilson coefficient is independent of the flow time to 
one loop and ${\cal O}(\tau)$, as we would expect:
\begin{align}
d_{\rho^2}\big\langle  \Omega 
|{} &\rho^2(\tau,0)\widetilde{\phi}(p_1)\widetilde{\phi}(p_2)
| \Omega \big\rangle =  1 \nonumber\\
{} & 
+\frac{\lambda}{32\pi^2}\left[2\gamma_{\;\mathrm{E}} + 
\log\left(\frac{m^2x^2}{4}\right)\right]+{\cal O}(\tau,\lambda^2).
\end{align}


\subsubsection{Next-to-leading-order}

To determine the next-to-leading-order contribution to the connected Wilson 
coefficient, we must incorporate the Feynman diagrams of 
Fig.~\ref{fig:sopehigher}. The two-loop diagrams of Fig.~\ref{fig:2loopexc} do 
not 
contribute to the 
Wilson coefficient $d_{\rho^2}$, because they appear at ${\cal O}(m^2)$.
\begin{figure}[htbp]\begin{minipage}{0.5\textwidth}
\includegraphics[width=0.95\textwidth,keepaspectratio=true]{images/dphi2loopm2}
\end{minipage}
\caption{\label{fig:2loopexc} 
Two-loop diagrams that appear at ${\cal O}(m^2)$ and therefore do not 
contribute to $d_{\rho^2}$. For details of the Feynman 
diagrams, see the caption of Fig.~\ref{fig:sope2pt}.}
\end{figure}

Looking at the Feynman diagrams of Fig.~\ref{fig:sopehigher}, and bearing 
in mind our experience of the next-to-leading-order contribution to 
$d_{\mathbb{I}}$, we immediately observe that these diagrams are simply the 
product of the one-loop Wilson coefficient, $d_{\rho^2}^{(1)}$, and the 
next-to-leading-order renormalized four-point vertex. This contribution to the 
vertex is, of course, nothing other than the next-to-leading-order contribution 
to the renormalized coupling constant. Thus, we write the Wilson 
coefficient quite simply as
\begin{equation}\label{eq:drho2r}
d_{\rho^2} = 1 +\frac{\lambda_{\mathrm{R}}}{32\pi^2}\left[\gamma_{\;\mathrm{E}} 
-1 + 
\log\left(\frac{x^2}{8\tau}\right)\right] +{\cal O}(\lambda^3),
\end{equation}
where $\lambda_{\mathrm{R}}$ is the renormalized coupling. In the $\ms$ scheme, 
the renormalized coupling is given by 
\begin{equation}\label{eq:lambdar}
\lambda_{\mathrm{R}} = \lambda - 
\frac{3\lambda^2}{2(16\pi^2)}\log\left(\frac{\mu^2}{m^2}\right)+{\cal 
O}(\lambda^3).
\end{equation}
For a calculation of this to five loops, see \cite{Kleinert:2001ax}.

As we move beyond the leading-order Wilson coefficients, 
\emph{i.e.}~$d_{\mathbb{I}}^{(0)}$ and $d_{\rho^2}^{(1)}$, divergent radiative 
corrections appear in our calculations, because all field interaction vertices 
appear at zero flow time. These divergences can be removed by the 
renormalization parameters of the original 
theory and the renormalization scale 
dependence of the smeared operators is completely contained in the 
renormalization parameters of the original theory. This is to be 
expected: it follows from the fact that renormalized matrix elements at zero 
flow time remain finite at nonzero flow time and require no further 
renormalization \cite{Makino:2014sta,Luscher:2011bx}.

Ultimately, for realistic calculations in lattice QCD, the 
perturbative 
calculation of smeared Wilson coefficients must be combined with 
nonperturbative determinations of matrix elements at hadronic energy 
scales. Scalar $\phi^4$ theory is not asymptotically free in four dimensions, 
but we can understand the mathematical features of the sOPE in more detail by 
studying the renormalization group equations for the simple example of 
scalar fields.
